% Part: set-theory
% Chapter: ordinals
% Section: wo

\documentclass[../../../include/open-logic-section]{subfiles}

\begin{document}

\olfileid{sth}{ordinals}{wo}
\olsection{সুক্রমবিন্যাস}

মৌলিক ধারণাটি এই:

\begin{defn}
$<$ সম্পর্কটি $A$-কে \emph{সুক্রমে বিন্যস্ত করে}
যদি এবং কেবল যদি নিচের দুটি শর্ত পূরণ হয়:
\begin{enumerate}
	\item $<$ সংযুক্ত; অর্থাৎ, সব $a, b \in A$-এর জন্য
	হয় $a < b$, নয় $a = b$, নয় $b < a$;
	\item $A$-এর প্রত্যেক অশূন্য উপসেটের একটি
	$<$-ন্যূনতম !!{element} আছে; অর্থাৎ,
	$\emptyset \neq X \subseteq A$ হলে $(\exists m \in
	X)(\forall z \in X)z \nless m$।
\end{enumerate}
% BN-SRC-784 TARGET-CORRECTION-BEGIN
\emph{পাঠ-স্পষ্টীকরণ।} এখানে ক্রম-সম্পর্ককে ক্ষেত্র $A$-তে
সীমিত করে নেওয়া হয়েছে। পরের আবেশে $b<a$-র পরিমাণায়নও
এই ক্ষেত্রেই; ক্ষেত্রের বাইরের বস্তুর উপর অতিরিক্ত দাবি নয়।
% BN-SRC-784 TARGET-CORRECTION-END
\end{defn}

এইমাত্র দেখা তিনটি দৃষ্টান্তই যে সুক্রমবিন্যাস,
তা সহজে যাচাই করা যায়।

\begin{prob}
\Olref[sth][ordinals][idea]{sec}-এ স্বাভাবিক সংখ্যার
উপর তিনটি ক্রম দেখানো হয়েছে। যাচাই করুন, প্রত্যেকটি
সুক্রমবিন্যাস।
\end{prob}

সুক্রমবিন্যাস সম্পর্কে কয়েকটি প্রাথমিক কিন্তু
অত্যন্ত গুরুত্বপূর্ণ কথা লক্ষ করি।

\begin{prop}\ollabel{wo:strictorder}
$<$ যদি $A$-কে সুক্রমে বিন্যস্ত করে, তবে $A$-এর
প্রত্যেক অশূন্য উপসেটের একটি একমাত্র $<$-ক্ষুদ্রতম
সদস্য আছে; আর $<$ আত্মসম্পর্কহীন, একমুখী ও পরিযায়ী।
\end{prop}

\begin{proof}
$X$ যদি $A$-এর অশূন্য উপসেট হয়, তবে এর একটি
$<$-ন্যূনতম !!{element} $m$ আছে; অর্থাৎ,
$(\forall z \in X)z \nless m$। যেহেতু $<$ সংযুক্ত,
তাই $(\forall z \in X)m \leq z$। অতএব $m$ হলো
$X$-এর $<$-ক্ষুদ্রতম !!{element}।

আত্মসম্পর্কহীনতা দেখাতে একটি $a \in A$ নিই।
$\{a\}$-এর $<$-ক্ষুদ্রতম !!{element} হলো $a$;
তাই $a \nless a$। পরিযায়িতা দেখাতে ধরি,
$a < b < c$। যেহেতু $\{a, b, c\}$-এর একটি
$<$-ক্ষুদ্রতম !!{element} আছে, তাই $a < c$।
আত্মসম্পর্কহীনতা ও পরিযায়িতা থেকে একমুখিতা পাওয়া যায়।
\end{proof}

\begin{prop}\ollabel{propwoinduction}
$<$ যদি $A$-কে সুক্রমে বিন্যস্ত করে, তবে যেকোনো
সূত্র $\phi(x)$-এর জন্য:
\[
	\text{যদি }(\forall a \in A)((\forall b < a)\phi(b) \lif
		\phi(a))\text{, তবে }(\forall a \in A)\phi(a).
\]
\end{prop}

\begin{proof}
আমরা প্রতিবিপরীত দাবিটি প্রমাণ করব। ধরি,
$\lnot(\forall a \in A)\phi(a)$; অর্থাৎ,
$X = \Setabs{x \in A}{\lnot\phi(x)} \neq \emptyset$।
তাহলে $X$-এর একটি $<$-ন্যূনতম !!{element} $a$ আছে।
অতএব $(\forall b < a)\phi(b)$, কিন্তু $\lnot \phi(a)$।
\end{proof}\noindent এই শেষ ধর্মটি স্বাভাবিক সংখ্যার উপর
শক্তিশালী আবেশের নীতির কথা মনে করায়। অর্থাৎ:
$(\forall n \in \omega)((\forall m < n)\phi(m)
\lif \phi(n))$ হলে $(\forall n \in \omega)\phi(n)$।
এই ধর্ম সুক্রমবিন্যাসকে খুবই \emph{শক্তিশালী}
ধারণা করে তোলে।\footnote{মনে রাখুন: আলাদা করে না
বললে সব সূত্রেই পরামিতি থাকতে পারে।}

\end{document}
